\chapter*{Abstract}
\phantomsection\addcontentsline{toc}{chapter}{\protect Abstract}
\thispagestyle{empty}

This thesis presents the development of an interferometric measurement system for the calibration of an automated translation stage used in a mid-infrared pump-probe spectroscopy experiment. The work was motivated by the limited positioning accuracy of the currently used control software. To improve the measurement of the stage displacement, a Michelson interferometer setup and corresponding analysis software were developed.

First, different laser sources were characterized and compared with respect to wavelength, beam quality, and intensity stability. Based on these measurements, the Thorlabs CPS780S laser was selected as the light source for the interferometer.

 Two methods for detecting the interference signal were compared: optical detection with a camera and electronic detection with a photodiode. Software was developed to evaluate both signals and count the interference fringes. The comparison showed that the camera-based detection was more stable and more reliable for the investigated setup.

Since a conventional Michelson interferometer does not provide information about the direction of motion, a quadrature homodyne interferometer was additionally developed. This extension was successfully built and showed that both displacement and direction of motion can be determined from two phase-shifted detector signals.

